\begin{abstract}
According to the United Nations Environment Programme (UNEP) \textit{Food Waste Index Report 2024}, approximately 1.05 billion tonnes of food were wasted globally across retail, food service, and household sectors in 2022, with households accounting for roughly 60\% of the total. Household meal preparation requires tracking ingredient quantities, portion scaling, date-label semantics, and food safety rules simultaneously; however, generic large language model (LLM) chatbots lack grounded inventory state and deterministic verification, often recommending infeasible dishes or omitting safety warnings. We present \textit{Food Guardian}, a hybrid AI agent integrating household inventory tracking, date-label reminders, preservation knowledge retrieval, inventory-aware meal recommendation, and consumption/waste analytics. Implemented with Streamlit and SQLite, Food Guardian orchestrates Gemini via LangGraph tool calling while delegating serving-scaled unit normalization, deficit calculation, dual-mode ranking (\texttt{ready\_now} and \texttt{expiry\_first}), and conservative safety pre-screening (\texttt{eligible}, \texttt{review}, \texttt{blocked}) to deterministic Python services. Preservation guidance combines structured JSON and a ChromaDB vector index with source attribution, and Pandas aggregates logged consumption and waste by compatible units. Under a specification-driven OpenSpec workflow (\texttt{FR-01}--\texttt{FR-10}), all 40 unit tests pass, 25 of 30 fixed evaluation cases pass offline deterministic/static checks (5 multi-tool cases remain unexecuted offline), and an initial Gemini Live run of 3 non-RAG multi-tool cases achieved 3/3 API success and 3/3 valid tool selection accuracy (11.71\,s mean latency). Planned paired comparisons against generic LLMs and live RAG groundedness audits are also formulated.
\end{abstract}

\begin{IEEEkeywords}
AI agent, household food waste, inventory-aware meal planning, large language models, tool calling, food safety screening, retrieval-augmented generation, OpenSpec.
\end{IEEEkeywords}
